%! Scatterplots and Correlation — Unit 2, Lesson 2.2 — Student Packet Cover Sheet
\documentclass[12pt]{article}
\usepackage{apstats-article}
\usepackage{apstats-boxes}
\usepackage{ltablex}
\keepXColumns

\begin{document}

% Full-bleed navy banner. \coverbanner MEASURES the title block and sizes the band to it.
\coverbanner{Unit 2}{Lesson 2.2 \quad Scatterplots and Correlation}

\namedateperiod
\vspace{0.18in}

\boxguard[14]
\begin{learningtargetbox}
\small
% GRADUAL RELEASE: the vocabulary is taught in the guided notes, so the targets name it
% plainly. The AP Learning Objectives (UNC-1.S, DAT-1.A, DAT-1.B, DAT-1.C) stay in the plan.
\begin{enumerate}[leftmargin=1.5em, itemsep=5pt]
  \item \ldots{} read a \textbf{scatterplot} and name its \textbf{explanatory variable} and its
        \textbf{response variable}.
  \item \ldots{} describe an \textbf{association} by its \textbf{direction}, \textbf{form},
        \textbf{strength}, and \textbf{unusual features} --- in context.
  \item \ldots{} interpret the \textbf{correlation} $r$ and explain why $r$ near 0 does not mean
        ``no relationship.''
  \item \ldots{} explain why a strong correlation does not show that one variable
        \textbf{causes} the other.
\end{enumerate}
\end{learningtargetbox}

\vspace{0.14in}

\boxguard[14]
\begin{tocbox}
\small
\renewcommand{\arraystretch}{1.5}
\begin{tabularx}{\linewidth}{@{} c l X r @{}}
\rowcolor{sky}
\textbf{\#} & \textbf{Component} & \textbf{Description} & \textbf{Score} \\
\midrule
1 & Warm-Up & Pass rates, the four words for a distribution, and five students' study
                 hours & \blank{1.2cm} \\
2 & Guided Notes \& Practice & Weight and fuel economy for 12 cars, one car driven at 12
                 speeds, and a city pool on hot days & \blank{1.2cm} \\
% AP Practice is assigned but NOT scored: its score cell prints NA, never a \blank{}.
% Row 4 is the lesson's GRADED practice. Paper homework is the DEFAULT for this course; on the
% occasions DeltaMath is assigned instead, that replaces row 4 and its score cell is struck.
3 & AP Practice & Screen time and sleep: two multiple-choice items and one free-response set
                 --- practice, not a grade & \textbf{NA} \\
4 & Homework & Practice hours and free throws, then what $r$ can and cannot tell you ---
                 \textbf{scored, due the first class after two study halls} & \blank{1.2cm} \\
\midrule
  & \multicolumn{2}{r}{\textbf{Total}} & \blank{1.2cm} \\
\end{tabularx}
\end{tocbox}

\vspace{0.14in}

\boxguard[8]
\begin{remindbox}
\small
\textbf{Look before you calculate.} The correlation $r$ only measures how well the points fit a
\emph{straight line}. A car driven at different speeds makes a clean curve and still has
$r$ close to 0 --- so an $r$ near 0 means ``no \emph{linear} association,'' not ``no
relationship.'' And a strong $r$ is not proof: two variables can rise together because
\textbf{something else} drives both. Describe what you see --- direction, form, strength,
unusual points --- in the words of the problem.
\end{remindbox}

\end{document}
